\begin{textsample}{POS Dim 1 – 2010\_c4 – Score -5.00 – 2010\_c4/enade\_2010\_tec\_agro\_ind\_8fdf27ddf8.txt}  \label{ex:f1_pos_401}
CAPÍTULO 4 PERCEPÇÃO DA PROVA As análises feitas neste capítulo tratam das percepções dos estudantes, separados por grupos de Ingressantes e Concluintes, sobre a prova aplicada no ENADE/2010 na Área de Tecnologia em Agroindústria.
Estas percepções foram mensuradas por meio de nove questões que avaliaram desde o grau de dificuldade da prova até o tempo gasto para concluí -la.
As percepções sobre a prova foram relacionadas com o desempenho dos alunos e com a Grande Região de funcionamento do curso.
O desempenho dos alunos foi classificado em dois níveis : quarto inferior e quarto superior.
Para tanto, esse desempenho foi colocado em ordem ascendente.
O percentil 25, P25, também conhecido como primeiro quartil, é a nota de desempenho que deixa um quarto ( 25,0 \% ) dos valores observados abaixo e três quartos acima.
A Figura 1 apresenta uma ilustração destes conceitos.
O quarto inferior de desempenho é composto pelas notas abaixo do primeiro quartil.
Já o percentil 75, P75, também conhecido como terceiro quartil, é o valor para o qual há três quartos ( 75,0 \% ) dos dados abaixo e um quarto acima dele.
O quarto superior de desempenho é composto pelas notas iguais ou acima do terceiro quartil.
O percentil 50, P50, também conhecido como mediana, é o valor que divide as notas em dois conjuntos de igual tamanho.
Vale ressaltar que percentis, quartis e medianas são pontos, e não obrigatoriamente pertencentes ao conjunto original de dados.
Ao passo que os quartos são subconjuntos dos dados originais.
A seguir, serão apresentados gráficos com resultados selecionados, relativos às nove questões avaliadas por grupos de estudantes.
As Tabelas no Anexo II apresentam os valores absolutos e a distribuição percentual das alternativas válidas das nove questões, segundo o mesmo recorte de desempenho dos alunos e Grande Região de funcionamento do curso. 4.1 GRAU DE DIFICULDADE DA PROVA 4.1.1 Componente de Formação Geral Na questão “ Qual o grau de dificuldade desta prova na parte de Formação Geral? ”, 20,6 \% do grupo de Ingressantes e 29,7 \% do grupo de Concluintes optaram pelas alternativas difícil ou muito difícil.
Entretanto, para mais da metade dos estudantes, tanto Ingressantes ( 65,1 \% ), quanto Concluintes ( 58,1 \% ) o Componente de Formação Geral da prova foi considerado com grau de dificuldade médio ( Gráfico 4.1 e Tabelas II.1 e II.10 no Anexo II ).
O percentual de Concluintes que consideraram a prova difícil ou muito difícil foi mais elevado no quarto inferior ( 43,8 \% ) do que no superior ( 21,1 \% ), o que também foi verificado no grupo dos Ingressantes, com percentual de participantes menor no quarto inferior ( 20,9 \% ), do que no quarto superior ( 27,1 \% ).
Semelhantemente ao ocorrido para os Ingressantes como um todo, para os quartos de menor e maior desempenho deste grupo, a alternativa modal para esta pergunta foi médio, com respectivamente 66,6 \% e 56,2 \% dos respondentes.
Já entre os Concluintes, estas percentagens foram, respectivamente, 49,9 \% e 63,1 \%.
Na prova de Tecnologia em Agroindústria, não houve a participação das cinco Grandes Regiões brasileiras para o grupo de Concluintes, tendo sido aplicada apenas nas regiões Norte e Sul do país para esses participantes.
Já para o grupo de Ingressantes, a prova foi aplicada em quatro das cinco Grandes Regiões, sendo que a Sudeste foi a única onde não houve ocorrência.
Para os Ingressantes, a região Nordeste foi aquela em que houve maior proporção dos que avaliaram o Componente de Formação Geral da prova como difícil ou muito difícil : 25,0 \% tiveram essa percepção da prova.
Dentre as duas regiões onde houve participantes Concluintes, a Sul foi aquela com percentual mais elevado dos que consideraram a parte de Formação Geral difícil ou muito difícil, com 53,1 \%, se comparada à Norte, com 11,9 \%.
Nas Grandes Regiões, a proporção de Ingressantes que consideraram o Componente de Formação Geral da prova como sendo de grau de dificuldade médio esteve localizada entre 62,4 \% na região Centro-Oeste e 66,5 \% na região Norte.
Nas duas Grandes Regiões participantes, no grupo de Concluintes, os percentuais estiveram entre 40,6 \% na Sul e 71,4 \% na região Norte. 4.1.2 Componente de Conhecimento Específico Ao responderem à questão “ Qual o grau de dificuldade desta prova na parte de Componente Específico? ”, 51,8 \% do grupo de Ingressantes e 43,2 \% do de Concluintes classificaram -na como difícil ou muito difícil.
O Componente de Conhecimento Específico da prova foi considerado com grau de dificuldade médio por 43,1 \% dos Ingressantes e por 51,3 \% dos Concluintes ( Gráfico 4.2 e Tabelas II.2 e II.11 no Anexo II ).
Considerando -se a percepção sobre o grau de dificuldade do Componente de Conhecimento Específico da prova de acordo com o desempenho dos estudantes, observa -se que 54,2 \% dos Ingressantes obtiveram desempenho situado no quarto superior classificaram -na como difícil ou muito difícil.
Dos Ingressantes que obtiveram desempenho situado no quarto inferior, 43,8 \% optaram pela mesma alternativa.
Para os Concluintes, o percentual de estudantes que consideraram o Componente de Conhecimento Específico da prova difícil ou muito difícil foi de 57,9 \% entre os que obtiveram desempenho no quarto superior e de 43,8 \% entre os situados no quarto inferior.
Para o quarto de melhor desempenho de Ingressantes, a alternativa modal para esta pergunta foi difícil, com 47,9 \% dos participantes.
Já entre os Concluintes, a alternativa modal para os do quadro superior também foi a opção difícil, com 52,6 \% dos respondentes.
Entre as Grandes Regiões brasileiras, a proporção de Ingressantes que avaliaram o Componente de Conhecimento Específico da prova como difícil ou muito difícil variou de 44,9 \% ( Sul ) a 66,7 \% ( Nordeste ), salientando -se que não houve a participação de Ingressantes na região Sudeste.
O grupo de Concluintes, com apenas duas regiões participantes, avaliou esse componente como difícil ou muito difícil, com 28,6 \% de participantes na região Norte e 62,6 \% na Sul.
O percentual de Ingressantes que classificaram o grau de dificuldade como médio, no Componente de Conhecimento Específico da prova, variou de 31,1 \% ( Centro-Oeste ) a 48,0 \% ( Sul ).
Dentre os Concluintes, a percepção do mesmo grau de dificuldade foi de 34,3 \%, na região Sul, e de 64,3 \%, na Norte. 4.2 Extensão da prova em relação ao tempo total Indagados quanto à extensão da prova, em relação ao tempo total oferecido para a sua resolução, os estudantes apontaram com maior incidência a alternativa que considerava a extensão adequada, para ambos os grupos e todas as agregações consideradas ( Gráfico 4.3 e Tabelas II.3 e II.12 no Anexo II ).
No grupo de Ingressantes, a maioria, 68,7 \%, respondeu ser a extensão da prova adequada, ao passo que, no grupo de Concluintes, este valor foi menor, 60,8 \%.
Para 29,2 \% dos Ingressantes o exame foi longo ou muito longo.
Essa mesma consideração foi feita por 31,1 \% dos Concluintes.
No grupo de Ingressantes, a maioria, tanto para o quarto de desempenho inferior ( 72,9 \% ), quanto para o de melhor desempenho ( 58,2 \% ) respondeu ser a extensão da prova adequada.
Já entre os Concluintes, esses valores foram menores, situados em cada quarto em, respectivamente, 56,1 \% e 57,9 \%.
Considerando -se ainda os quartos extremos de desempenho, observa -se que a proporção de estudantes que consideraram a prova longa ou muito longa, em relação ao tempo total destinado à sua resolução, acompanhou a percepção de Ingressantes e Concluintes.
Para o grupo de Ingressantes, entre os estudantes no quarto de desempenho inferior, 27,1 \% declararam ter esta percepção, ao passo que, entre aqueles no quarto de melhor desempenho, esta proporção foi 35,5 \%.
O mesmo pode ser constatado entre os Concluintes : 25,1 \% no quarto de menor desempenho e 36,8 \% no quarto de maior desempenho.
Na análise por Grandes Regiões, o percentual de Ingressantes que responderam ser a extensão da prova adequada variou de 62,3 \% no Sul a 83,3 \% no Nordeste.
Já entre os Concluintes, estes valores, nas duas regiões participantes, foram de 59,4 \% na Sul a de 62,0 \% na Norte.
Aqueles Ingressantes que avaliaram a prova como longa ou muito longa em relação ao tempo total destinado à sua resolução variou de 16,7 \% na região Nordeste a 35,7 \% na região Sul.
Os percentuais quanto ao mesmo aspecto, entre os Concluintes, nas regiões participantes, foram de 28,5 \% na região Norte a 34,4 \% na região Sul. 4.3 COMPREENSÃO DOS ENUNCIADOS DAS QUESTÕES 4.3.1 Componente de Formação Geral Com relação aos enunciados das questões do Componente de Formação Geral, as opiniões foram positivas, com percentuais que foram dos 74,3 \% de Concluintes aos 82,0 \% de Ingressantes que consideraram claros e objetivos os enunciados de todas ou da maioria das questões ( Gráfico 4.4 e Tabelas II.4 e II.13 no Anexo II ).
Ao se comparar as percepções dos estudantes segundo o desempenho, observou -se que, entre os Ingressantes do quarto superior, a clareza e objetividade de todos ou da maioria dos enunciados das questões foram percebidas por 81,2 \%.
A proporção apurada para os Concluintes desse mesmo quarto foi 89,5 \%.
No quarto de desempenho inferior, tal avaliação foi emitida por 79,1 \% dos Ingressantes e por 74,9 \% dos Concluintes.
A partir dessas estatísticas, é possível afirmar que todos ou a maioria dos enunciados das questões do Componente de Formação Geral estavam claros e objetivos para a maioria dos respondentes, segundo os quartos de desempenho considerados, embora os que se situaram no quarto superior tenham declarado esta percepção com maior incidência.
Na análise regional, observa -se que as opiniões de que todos ou a maioria dos enunciados das questões do Componente de Formação Geral estavam claros e objetivos variaram, entre os Ingressantes, de 62,4 \% no Centro-Oeste a 87,5 \% no Nordeste.
Entre os Concluintes, os percentuais foram 71,5 \% no Norte e 78,1 \% no Sul, sendo estas as duas regiões participantes. 4.3.2 Componente de Conhecimento Específico Com relação aos enunciados das questões do Componente de Conhecimento Específico da prova, para os Concluintes da Área de Tecnologia em Agroindústria, a clareza e a objetividade foram percebidas em todos ou na maioria destes por 75,4 \%.
Para o grupo dos Ingressantes, foram percebidas por 72,7 \%.
A proporção de estudantes que considerou os enunciados das questões claros e objetivos foi, nos quartos superiores, de 77,0 \% dos Ingressantes e 78,9 \% dos Concluintes.
Nos quartos inferiores, foram de 78,8 \% dos Ingressantes e de 66,7 \% dos Concluintes ( Gráfico 4.5 e Tabelas II.5 e II.14 no Anexo II ).
Nas duas Grandes Regiões brasileiras participantes, o percentual de Concluintes que consideraram claros e objetivos todos ou a maioria dos enunciados das questões do Componente de Conhecimento Específico da prova foi de 68,7 \% no Sul e de 80,5 \% no Norte.
Os Ingressantes da região Nordeste ( 66,6 \% ) apresentaram proporção mais baixas, enquanto os da região Sul ( 76,3 \% ), a mais elevada desta percepção. 4.4 SUFICIÊNCIA DAS INFORMAÇÕES/INSTRUÇÕES FORNECIDAS Ao avaliarem as informações/instruções fornecidas para a resolução das questões, 82,1 \% do grupo dos Ingressantes e 83,7 \% do grupo dos Concluintes respondentes de todo o Brasil afirmaram que estas eram até excessivas ou suficientes em todas ou na maioria das questões ( Gráfico 4.6 e Tabelas II.6 e II.15 no Anexo II ).
Levando -se em conta o desempenho dos participantes, notam -se diferenças entre as opiniões de estudantes, Ingressantes e Concluintes, dos quartos inferior e superior de desempenho, como mostra o Gráfico 4.6.
O percentual de participantes que avaliaram as informações/instruções como excessivas ou suficientes em todas ou na maioria das questões foi mais elevado nos quartos superiores, correspondendo a 87,5 \% dos Ingressantes e a 89,5 \% dos Concluintes, percentuais superiores à média nacional para cada um dos dois grupos.
Já nos quartos inferiores, a suficiência das informações/instruções declarada como excessiva em todas ou na maioria das questões foi percebida com o mesmo percentual, por 81,2 \% dos Ingressantes e dos Concluintes.
Quanto à distribuição de respondentes nas Grandes Regiões, observa -se que as proporções de estudantes que perceberam como excessivas ou suficientes em todas ou na maioria das questões as informações/instruções fornecidas para a resolução das mesmas superaram as médias nacionais, em maior medida, entre os Ingressantes da região SNordeste ( 91,7 \% ).
Entre os Concluintes, os percentuais para as duas regiões participantes foram 81,2 \% na Sul e 85,7 \% na Norte. 4.5 DIFICULDADE ENCONTRADA AO RESPONDER À PROVA Perguntados sobre as dificuldades com as quais se depararam ao responder à prova, 57,6 \% do grupo de Ingressantes e 12,2 \% do grupo de Concluintes apontaram o desconhecimento do conteúdo.
Concluintes, por sua vez, apontaram a forma diferente de abordagem do conteúdo como dificuldade, mencionada por 39,2 \%.
Esta alternativa foi apontada por uma parcela menor dos Ingressantes, 25,3 \%.
A falta de motivação para fazer a prova foi a dificuldade para 24,3 \% dos Concluintes e para 9,3 \% dos Ingressantes.
Considerando -se os respondentes em todo o Brasil, 5,7 \% dos Ingressantes e 16,2 \% dos Concluintes afirmaram que não tiveram qualquer tipo de dificuldade para responder à prova ( Tabelas II.7 e II.16 no Anexo II ).
Com relação aos quartos de desempenho, o desconhecimento do conteúdo foi a dificuldade modal entre os Ingressantes, apontada por 65,9 \% daqueles no quarto superior e por 58,3 \% dos do inferior.
Entre os Concluintes, os percentuais de estudantes com a mesma percepção foi inferior ao encontrado entre os Ingressantes, 15,8 \% no quarto superior e 0,0 \% no quarto inferior.
A alternativa modal no grupo de Concluintes, quando desagregada pelo desempenho, foi que a dificuldade encontrada era a forma diferente de abordagem do conteúdo : 43,7 \% no quarto inferior e 47,4 \% do quarto superior assim o responderam.
Na análise por Grandes Regiões, o percentual de Ingressantes que apontaram o desconhecimento do conteúdo como dificuldade ao responder à prova variou de 51,6 \%, na região Sul a 79,2 \%, na Nordeste, recordando a não participação da região Sudeste.
Entre os Concluintes, o desconhecimento do conteúdo foi a resposta de 11,9 \% dos participantes da região Norte e de 12,5 \% da região Sul.
A forma diferente de abordagem do conteúdo foi citada, entre os Concluintes, por 31,0 \% ( região Norte ) e 50,0 \% ( Sul ).
Os percentuais relacionados a essa opção entre os Ingressantes variaram de 12,5 \% ( região Nordeste ) a 30,9 \% ( região Sul ).
O percentual de Ingressantes que citou a falta de motivação como dificuldade superou a média nacional na região Norte ( 15,8 \% ).
Entre os Concluintes, os percentuais para essa opção foram 9,4 \% ( região Sul ) e 35,7 \% ( região Norte ).
Houve percentuais acima da média nacional dos Ingressantes nas regiões Sul ( 7,2 \% ) e Nordeste ( 8,3 \% ), com declarações de que não tiveram qualquer dificuldade para responder à prova.
Foram 9,5 \% dos Concluintes da região Norte e 25,0 \% dos da região Sul a emitirem a mesma opinião.
O Gráfico 4.7 apresenta os percentuais de estudantes, Ingressantes e Concluintes, que apontaram o desconhecimento do conteúdo como dificuldade percebida ao responder à prova. 4.6 CONTEÚDOS DAS QUESTÕES OBJETIVAS DA PROVA Ao analisarem os conteúdos das questões objetivas da prova, a maior parte dos Ingressantes, 70,1 \% do total nacional, afirmou que não estudou ainda a maioria desses conteúdos, enquanto entre os Concluintes, apenas 2,7 \% expressaram a mesma opinião.
Por outro lado, a maioria dos Concluintes ( 74,3 \% ) afirmou ter estudado e aprendido muitos ou todos os conteúdos.
Para os Ingressantes, o percentual correspondente a essas alternativas foi 18,5 \% ( Gráfico 4.8 e Tabelas II.8 e II.17 no Anexo II ).
Considerando -se as opiniões de estudantes de acordo com os quartos de desempenho, os Ingressantes do quarto superior ( 68,1 \% ) obtiveram percentual semelhante ao dos Ingressantes do quarto inferior ( 68,8 \% ), ao afirmarem que não estudaram ainda a maioria desses conteúdos.
Para os Concluintes do quarto inferior, 6,3 \% ofereceram essa resposta, sendo que 0,0 \% dos Concluintes emitiu essa opinião no quarto superior.
Tendo em conta os quartos superiores para os dois grupos, 84,2 \% dos Concluintes e 29,8 \% dos Ingressantes afirmaram ter estudado e aprendido muitos ou todos os conteúdos.
Na análise por Grande Região, a proporção de Ingressantes que afirmou que não estudou ainda a maioria desses conteúdos superou a média nacional nas regiões Norte ( 80,6 \% ) e Nordeste ( 83,4 \% ).
Entre os Concluintes, nesta mesma alternativa, houve menção apenas na região Sul, com 6,3 \% de respondentes.
Também entre os Concluintes, 88,1 \% dos da região Norte e 56,2 \% da região Sul afirmaram ter estudado e aprendido muitos ou todos os conteúdos. 4.7 TEMPO GASTO PARA CONCLUIR A PROVA Ao responderem sobre o tempo de conclusão da prova, mais da metade dos estudantes afirmou ter gasto entre duas e quatro horas.
Considerando -se separadamente as percepções do grupo de Ingressantes e do grupo de Concluintes, os percentuais, em termos nacionais, foram de 71,9 \% e 70,3 \%, respectivamente ( Gráfico 4.9 e Tabelas II.9 e II.18 no Anexo II ).
Uma vez sendo analisadas as alternativas oferecidas pelos estudantes que se situaram nos quartos superior e inferior, observa -se que uma maior proporção de Ingressantes no quarto superior declarou ter gasto entre duas e quatro horas para concluir a prova quando comparados com os do quarto inferior, respectivamente 77,1 \% e 54,1 \%.
A mesma situação foi verificada no grupo dos Concluintes, no qual 84,2 \% daqueles no quarto superior e 43,7 \% daqueles no quarto inferior declaram ter gasto entre duas e quatro horas para concluir a prova.
Nas cinco Grandes Regiões brasileiras, os Ingressantes que utilizaram entre duas e quatro horas para finalizar a prova superou o percentual nacional, em maior medida, a região Nordeste ( 82,7 \% ).
Entre os Concluintes, dispensaram entre duas e quatro horas para concluir a prova 46,8 \% da região Sul e 88,1 \% da Norte, como mostra o Gráfico 4.9.

% matched lemmas: 
\end{textsample}
